\documentclass[10pt,a4paper]{amsart}
\usepackage[margin=19mm]{geometry}
\usepackage{amsmath,amssymb,mathtools}
\usepackage[scr=rsfs,cal=euler]{mathalfa}
\usepackage{xfrac}
\usepackage[unicode,hidelinks,bookmarks=false]{hyperref}
\newcommand{\ie}{i.e.\ }
\newcommand{\cf}{cf.\ }
\newcommand{\resp}{respectively\ }
\newcommand{\Subsection}{\subsection}
\providecommand{\nobreakdash}{\nobreak\hbox{-}}
\setlength{\emergencystretch}{2em}
\multlinegap0pt
\usepackage{bm}
\usepackage{bbm}
\usepackage{dsfont}
\usepackage{xspace}
\usepackage{cancel}
\usepackage{mathtools}
\usepackage{amsthm}
\makeatletter
\@ifundefined{lemma}{\newtheorem{lemma}{Lemma}}{}
\makeatother
\title[Dynamical low-rank approximation for the semiclassical Schro]{Dynamical low-rank approximation for the semiclassical Schrodinger equation with uncertainties}
\author{Liu Liu, Limin Xu, Zhenyi Zhu}
\date{Source: arXiv:2602.06808v2 (CC BY 4.0)}
\begin{document}
\maketitle

\noindent\emph{Source and licence.} Dynamical low-rank approximation for the semiclassical Schrodinger equation with uncertainties. Version arXiv:2602.06808v2 (CC BY 4.0), distributed under the Creative Commons Attribution 4.0 International licence, as recorded in the arXiv licence metadata for this exact version and shown on the abstract page https://arxiv.org/abs/2602.06808v2. Full rights notice: licenses/arxiv-2602.06808-CC-BY-4.0.txt.\par\smallskip

\noindent\textit{Editorial scope.} Counted unit: the Lemma whose source label is \texttt{None} in the article (source0.tex lines 363--379, source label \texttt{None}), delivered together with the article's own complete original proof of it (source0.tex lines 380--427, one \texttt{proof} environment of 48 lines). No mathematics is written, repaired or re-derived: statement and proof are the article's own text line for line. Required context retained uncounted: none. Every definition, display and label of those lines is retained unchanged. Omitted from this package and not redistributed: 1--362, 428--1903. Nothing inside a retained excerpt is omitted or altered: every retained body is delivered exactly as the article's lines are written, so no omission receipt of any kind accompanies this package. Citations: the retained excerpts carry no citation command at all, so no bibliography entry of the article is transcribed and no reference is redistributed. Reference adaptation: none was needed and none was made; every reference inside the delivered unit resolves inside it, so no unresolved reference or citation is produced. Printed slips kept literal: no printed word, symbol, hypothesis or number of the counted statement or of its proof was altered. Layout and class: the article's own class is \texttt{elsarticle} with packages including geometry, amssymb, amsthm, amsmath, amsfonts, hyperref, cleveref, graphicx, mathtools, enumitem, listings, bm, color, xcolor; the wrapper is the lane's pinned \texttt{amsart} editorial wrapper at 10pt on A4, which restates the theorem-like environments the retained text needs and adds the article's own macro definitions that the retained text needs (none), with extra wrapper packages (bm, bbm, dsfont, xspace, cancel, mathtools). The delivered numbering is the unit's own. Assets: no retained span contains a figure environment, a graphics-inclusion command, a PDF-page inclusion, a table or a TikZ picture; no image file is copied into this package, the package declares no asset dependency and no image file is redistributed with it. Licence: version arXiv:2602.06808v2 of this article is distributed under the Creative Commons Attribution 4.0 International licence, as recorded in the arXiv licence metadata for this exact version and shown on the abstract page https://arxiv.org/abs/2602.06808v2. A full rights notice is delivered at licenses/arxiv-2602.06808-CC-BY-4.0.txt. Characters: every retained excerpt is pure ASCII, so the delivered engine, XeLaTeX with its default Latin Modern text font, reports no missing character.
\begin{lemma}
    The solutions $\psi^{(k)}(t) = \sum_{i,j=1}^r X_i^{(k)}(t) S_{ij}^{(k)}(t) W_j^{(k)}(t)$ for $k=1,2,3$ are determined as follows:
    \begin{enumerate}
        \item \textbf{K-step}: $W_j^{(1)}$ remains constant ($W_j^{(1)}(t) = W_j(t_n)$). The updated spatial factors $K_j(t, x) = \sum_{i=1}^r X_i^{(1)}(t) S_{ij}^{(1)}(t)$ evolve according to:
        \begin{align*}
            \partial_t K_j(t, x) = \langle F(t, \sum_{j=1}^r K_j(t, x) W_j(t_n, \xi)), W_j(t_n, \xi) \rangle_{\xi}, \quad j=1,\dots,r.
        \end{align*}
        \item \textbf{S-step}: Both $X_i^{(2)}$ and $W_j^{(2)}$ remain constant ($X_i^{(2)} = X_i^{(1)}(t_{n+1}), W_j^{(2)} = W_j(t_n)$). The core matrix $S_{ij}^{(2)}$ evolves according to:
        \begin{align*}
            \partial_t S_{ij}^{(2)} = -\langle F(t, \sum_{k,l=1}^r X_k^{(2)} S_{kl}^{(2)} W_l^{(2)}), X_i^{(2)} W_j^{(2)} \rangle_{x,\xi}, \quad i,j=1,\dots,r.
        \end{align*}
        \item \textbf{L-step}: $X_i^{(3)}$ remains constant ($X_i^{(3)} = X_i^{(1)}(t_{n+1})$). The updated random factors $L_i(t, \xi) = \sum_{j=1}^r S_{ij}^{(3)}(t) W_j^{(3)}(t)$ evolve according to:
        \begin{align*}
            \partial_t L_i(t, \xi) = \langle F(t, \sum_{i=1}^r X_i^{(3)}(x) L_i(t, \xi)), X_i^{(3)}(x) \rangle_{x}, \quad j=1,\dots,r.
        \end{align*}
    \end{enumerate}
\end{lemma}

\begin{proof}
    \textit{Step 1: The K-step.} 
    Let $\psi^{(1)} = \sum_{j=1}^r K_j^{(1)} W_j^{(1)}$, where the spatial factors are defined as $K_j^{(1)} = \sum_{i=1}^r X_i^{(1)} S_{ij}^{(1)}$. The evolution equation for this step is given by:
    \begin{align*}
        \partial_t \psi^{(1)} = P_{\mathcal{V}_W} F(\psi^{(1)}) = \sum_{j=1}^r \langle F(\psi^{(1)}), W_j^{(1)} \rangle_{\xi} W_j^{(1)}.
    \end{align*} 
    Substituting the expression for $\psi^{(1)}$ into the above yields:
    \begin{align}\label{eq: proof_step1_expand}
        \sum_{j=1}^r \left( \partial_t K_j^{(1)} W_j^{(1)} + K_j^{(1)} \partial_t W_j^{(1)} \right) = \sum_{j=1}^r \langle F(\psi^{(1)}), W_j^{(1)} \rangle_{\xi} W_j^{(1)}.
    \end{align}
    Taking the inner product with $W_k^{(1)}$ on both sides, and invoking the orthonormal property $\langle W_j^{(1)}, W_k^{(1)} \rangle_{\xi} = \delta_{jk}$ along with the gauge condition $\langle \partial_t W_j^{(1)}, W_k^{(1)} \rangle_{\xi} = 0$, we obtain:
    \begin{align*}
        \partial_t K_k^{(1)} = \langle F(\psi^{(1)}), W_k^{(1)} \rangle_{\xi}.
    \end{align*}
    Substituting this back into \eqref{eq: proof_step1_expand} implies $\sum_{j=1}^r K_j^{(1)} \partial_t W_j^{(1)} = 0$. By writing $K_j^{(1)} = \sum_i X_i^{(1)} S_{ij}^{(1)}$ and taking the inner product with $X_i^{(1)}$, the orthonormality of $X_i^{(1)}$ and the nonsingularity of $S^{(1)}$ lead to $\partial_t W_j^{(1)} = 0$.

    \textit{Step 2: The S-step.} 
    For $\psi^{(2)} = \sum_{i,j=1}^r X_i^{(2)} S_{ij}^{(2)} W_j^{(2)}$, the evolution is governed by:
    \begin{align*}
        \partial_t \psi^{(2)} = -P_{\mathcal{V}_W} P_{\mathcal{V}_X} F(\psi^{(2)}) = -\sum_{i,j=1}^r \langle F(\psi^{(2)}), X_i^{(2)} W_j^{(2)} \rangle_{x,\xi} X_i^{(2)} W_j^{(2)}.
    \end{align*}
    Expanding the time derivative $\partial_t \psi^{(2)}$ gives:
    \begin{equation}\label{eq: proof_step2_expand}
    \begin{aligned}
        &\sum_{i,j=1}^r \left( \partial_t X_i^{(2)} S_{ij}^{(2)} W_j^{(2)} + X_i^{(2)} \partial_t S_{ij}^{(2)} W_j^{(2)} + X_i^{(2)} S_{ij}^{(2)} \partial_t W_j^{(2)} \right) \\
        =& -\sum_{i,j=1}^r \langle F(\psi^{(2)}), X_i^{(2)} W_j^{(2)} \rangle_{x,\xi} X_i^{(2)} W_j^{(2)}.
    \end{aligned}
    \end{equation}
    % \begin{equation}\label{eq: proof_step2_expand}
    %     \sum_{i,j} \left( \partial_t X_i^{(2)} S_{ij}^{(2)} W_j^{(2)} + X_i^{(2)} \partial_t S_{ij}^{(2)} W_j^{(2)} + X_i^{(2)} S_{ij}^{(2)} \partial_t W_j^{(2)} \right) = \partial_t \psi^{(2)}.
    % \end{equation}
    Taking the inner product with $X_k^{(2)} W_l^{(2)}$ and applying the gauge conditions $\langle \partial_t X_i, X_k \rangle_x = 0$ and $\langle \partial_t W_j, W_l \rangle_\xi = 0$, we isolate the core matrix evolution:
    \begin{align*}
        \partial_t S_{ij}^{(2)} = -\langle F(\psi^{(2)}), X_i^{(2)} W_j^{(2)} \rangle_{x,\xi}.
    \end{align*}
    Substituting this result into \eqref{eq: proof_step2_expand} yields $\sum_{i,j} (\partial_t X_i^{(2)} S_{ij}^{(2)} W_j^{(2)} + X_i^{(2)} S_{ij}^{(2)} \partial_t W_j^{(2)}) = 0$. Using the orthonormality and nonsingularity arguments as in Step 1, we conclude $\partial_t X_i^{(2)} = 0$ and $\partial_t W_j^{(2)} = 0$.

    \textit{Step 3: The L-step.} 
    Let $\psi^{(3)} = \sum_{i=1}^r X_i^{(3)} L_i^{(3)}$ with $L_i^{(3)} = \sum_{j=1}^r S_{ij}^{(3)} W_j^{(3)}$. The governing equation is:
    \begin{align*}
        \partial_t \psi^{(3)} = P_{\mathcal{V}_X} F(\psi^{(3)}) = \sum_{i=1}^r \langle F(\psi^{(3)}), X_i^{(3)} \rangle_{x} X_i^{(3)}.
    \end{align*}
    Repeating the projection procedure by taking the inner product with $X_k^{(3)}$, we find:
    \begin{align*}
        \partial_t L_i^{(3)} = \langle F(\psi^{(3)}), X_i^{(3)} \rangle_{x}.
    \end{align*}
    Similarly, substituting this back into the evolution of $\psi^{(3)}$ and utilizing the orthonormality of $W_j$ and nonsingularity of $S^{(3)}$, we obtain $\partial_t X_i^{(3)} = 0$, which concludes the proof.
\end{proof}

\end{document}
